% title:          Sinking the polynomial submarine
% area:           strategy-puzzles, algorithms
% methods:        bijection, construction
% difficulty:
% difficulty-ai:  easy
% status:
% review:         none
% --- history: all optional, leave EMPTY if unknown ---
% origin:         folklore
% origin-date:    2002
% origin-number:
% also-in:        course
% taken-from:
% references:     Linear version (integer start, integer velocity): wu::riddles (William Wu, UC Berkeley), hard riddles, 'Sink The Sub', contributed by [srowen], 2002 - https://web.archive.org/web/20021007223120/http://www.ocf.berkeley.edu/~wwu/riddles/hard.shtml; wu::forums thread 'Sink The Sub', opened 31 Jul 2002 (reply #18 by W. Wu, 9 Apr 2004: any finite number of countably-valued parameters still works; reply #19, 29 Apr 2004, suggests a Taylor-polynomial generalisation) - https://web.archive.org/web/20250818025840/https://www.ocf.berkeley.edu/~wwu/cgi-bin/yabb/YaBB.cgi?board=riddles_hard;action=display;num=1028142186; R. Vakil, Stanford Problem-Solving Masterclass, Week 7 (29 Nov 2004), Problem 4, credited to Wu's website - http://math.stanford.edu/~vakil/putnam04/04mc7.pdf; Math Fail, 'Submarine Puzzle', Jan 2011 - https://www.math-fail.com/2011/01/submarine-puzzle.html; 'Catching a spy' variant (yes/no queries), K. R. M. Leino, from Radu Grigore - https://leino.science/puzzles/catching-a-spy; real-parameter variant with a bomb of width 1: Math StackExchange Q642957 (2014) - https://math.stackexchange.com/questions/642957. No published source found for the rational-polynomial, unknown-degree version; it looks like the club's own generalisation.
% history:        ai
\begin{problem}
An enemy submarine is hidden somewhere along the infinite line $\mathbb{R}$. It travels silently, and you know that its path is described by a fixed rational polynomial rule \(\sum_{i=0}^{n}c_iT^{i}\in\mathbb{Q}[X]\); for each time \(t \in \mathbb{N}\), its position is given by
\[
x(t) = \sum_{i=0}^{n} c_i t^{i}\in\mathbb{Q}.
\]
However, you do not know \(n \in \mathbb{N}\), nor the rational coefficients \((c_i)_{i \in n} \in \mathbb{Q}^n\) that form this rational polynomial rule.
\\\\
Each unit of time (starting from \(0\)), you are allowed to launch a single torpedo at any chosen rational position. If the submarine is at that position at that time, it is struck and sinks. Assume you possess a potentially countably infinite arsenal of torpedoes and potentially countably infinite time.
\\\\
Devise a strategy — a sequence of torpedo launches — such that, regardless of the submarine's unknown position, you will \textbf{eventually} hit it.
\end{problem}
\begin{solution}
\[
\substack{\textit{This solution needs to be written more formally}}
\]
Notice that \(\forall t \in \mathbb{N}\), we have \(x(t) \in \mathbb{Q}\). We must find a function
\[
f: \mathbb{N} \rightarrow \mathbb{Q}
\]
(which represents the sequence of torpedo launches — one per unit of time) such that there exists a time \(\tilde{t} \in \mathbb{N}\) with
\[
f(\tilde{t}) = x(\tilde{t}).
\]
The strategy is the following: assume the submarine does not move, i.e. \(x(T) = c_0 \in \mathbb{Q}[T]\), then taking any bijection \(g: \mathbb{N} \simeq \mathbb{Q}\) and defining \(f(t) = g(t)\) works.
\\\\
If the submarine moves linearly, \(x(T) = c_0 + c_1 T \in \mathbb{Q}[T]\), since we do not know \(c_0, c_1\), we enumerate each 2-tuple of rationals in the following way: suppose it started with \((c_0', c_1')\), and knowing the elapsed time \(t\), we shoot at position \(c_0' + c_1' t\). If it truly started with \((c_0', c_1')\), then we would certainly hit it. If it did not start with \((c_0', c_1')\) and by chance we still hit it, we are done; otherwise, we proceed to try another pair \((c_0'', c_1'')\), and continue this process, incrementing time to \(t' = t + 1\) each unit of time. Taking any bijection \(g: \mathbb{N} \simeq \mathbb{Q}^{2}\) and defining \(f(t) = g(t)(0)+g(t)(1)t\) works.
\\\\
If the submarine moves quadratically, \(x(T) = c_0 + c_1 T + c_2 T^2 \in \mathbb{Q}[T]\), again, since we do not know \(c_0, c_1, c_2\), we enumerate each 3-tuple of rationals similarly: suppose it started with \((c_0', c_1', c_2')\), and knowing the elapsed time \(t\), we shoot at position \(c_0' + c_1' t + c_2' t^2\). If it truly started with \((c_0', c_1', c_2')\), we would certainly hit it. If it did not start with that and we still hit it by chance, we are done; otherwise, we proceed to try another triple \((c_0'', c_1'', c_2'')\), incrementing time each step. Taking any bijection \(g: \mathbb{N} \simeq \mathbb{Q}^{3}\) and defining \(f(t) = g(t)(0)+g(t)(1)t+g(t)(2)t^2\) works.
\\\\
The problem you may raise is: we already enumerate \(\mathbb{Q}\) in the case the submarine does not move. How can we have finished that enumeration and moved on to the linear case, and further to the quadratic case, etc.?
\\\\
The key observation is that we can enumerate all of them simultaneously in a single shot! Indeed, we can enumerate all \textit{finite} sequences of rationals because
\[
\mathbb{N} \simeq \bigcup_{k \in \mathbb{N}} \mathbb{Q}^k =: \mathbb{Q}^{<\omega}.
\]
Any bijection \(g: \mathbb{N} \simeq \mathbb{Q}^{<\omega}\) suffices, since then the strategy at time \(t \in \mathbb{N}\) is, if \(g(t) = (c_i')_{i \in k}\), to shoot at \(\sum_{i \in k} c_i' t^{i}\). That is, the sequence of torpedo launches \(f: \mathbb{N} \rightarrow \mathbb{Q}\) is defined by:
\[
f(t) := \sum_{i \in \operatorname{dom}(g(t))} g(t)(i) t^{i}.
\]
You may object that such a bijection \(g: \mathbb{N} \rightarrow \mathbb{Q}^{<\omega}\) requires the axiom of choice and is not realisable in practice; however, we prove otherwise by providing multiple constructive examples of such bijections.
\\\\
First, we show that \(\mathbb{N} \simeq \mathbb{Q}\). Define \(s: \mathbb{N} \to \mathbb{Z}\) for \(n \in \mathbb{N}\) by:
\[
s(n) = (-1)^n \left\lfloor \frac{n + 1}{2} \right\rfloor,
\]
which is clearly a bijection. (The sequence \(s(0), s(1), s(2), s(3), \ldots\) is \(0, -1, 1, -2, 2, \ldots\), a bijection from \(\mathbb{N}\) to \(\mathbb{Z}\).)
\\\\
It suffices then to find a bijection between \(\mathbb{Z}\) and \(\mathbb{Q}\). Any bijection \(h_+: \mathbb{Z}_{>0} \to \mathbb{Q}_{>0}\) induces a bijection \(h: \mathbb{Z} \to \mathbb{Q}\) defined by:
\[
h(n) =
\begin{cases}
h_+(n) & \text{if } n > 0, \\
0 & \text{if } n = 0, \\
-h_+(-n) & \text{if } n < 0.
\end{cases}
\]
It remains to find a bijection \(h_+: \mathbb{Z}_{>0} \to \mathbb{Q}_{>0}\). Every positive integer has a unique prime factorisation \(p_1^{a_1} p_2^{a_2} \cdots\), where \(p_i\) is the \(i\)-th prime and the \(a_i\) are non-negative integers, almost all zero. Similarly, every positive rational number has a unique expression \(p_1^{b_1} p_2^{b_2} \cdots\) with finitely many non-zero integers \(b_i\).
\\\\
Thus, for any \(n = p_1^{a_1} p_2^{a_2} \cdots \in \mathbb{Z}_{>0}\), define
\[
h_+(n) := p_1^{s(a_1)} p_2^{s(a_2)} \cdots.
\]
This defines a bijection \(h_+: \mathbb{Z}_{>0} \to \mathbb{Q}_{>0}\), and \(h\) gives a bijection \(\mathbb{Z} \to \mathbb{Q}\), hence \(\mathbb{N} \simeq \mathbb{Q}\).
\\\\
Such a bijection induces one between \(\mathbb{N}^{<\omega}\) and \(\mathbb{Q}^{<\omega}\). So, to get a bijection between \(\mathbb{N}\) and \(\mathbb{Q}^{<\omega}\), it suffices to give one between \(\mathbb{N}\) and \(\mathbb{N}^{<\omega}\). We give three examples \(f_1, f_2, f_3: \mathbb{N}^{<\omega} \rightarrow \mathbb{N}\) below:
\\\\
The prime encoding bijection:
\[
f_1((a_0, a_1, \dots, a_{k-1})) := \prod_{i=0}^{k-1} p_i^{a_i + 1} - 1,
\]
where \(p_i\) is the \(i\)-th prime. Note \(f_1(\varepsilon) = 0\).
\\\\
The multiset rank encoding bijection:
\[
f_2((a_0, a_1, \dots, a_{k-1})) := \sum_{j=0}^{k-1} \binom{a_j + j}{j + 1}.
\]
Again, \(f_2(\varepsilon) = 0\).
\\\\
The recursive Cantor pairing bijection:
\\\\
Recall the \href{https://en.wikipedia.org/wiki/Pairing_function}{Cantor pairing function}:
\[
\pi(a, b) := \frac{1}{2}(a + b)(a + b + 1) + b.
\]
Define recursively:
\[
f_3(\varepsilon) := 0, \quad f_3((a_0, \dots, a_k)) := \pi(a_0, f_3((a_1, \dots, a_k))) + 1.
\]
These induce constructive bijections \(g_1, g_2, g_3: \mathbb{N} \to \mathbb{Q}^{<\omega}\).
\\\\
Now, we show that any bijection \(g: \mathbb{N} \to \mathbb{Q}^{<\omega}\) induces a well-defined strategy that works. Technically, we only use the surjectivity of \(g\). Let \(n\) and \(\mathbf{c} := (c_i)_{i \in n} \in \mathbb{Q}^n\) be the true coefficients of the submarine's trajectory. Then by surjectivity of \(g\), there exists \(\tilde{t} \in \mathbb{N}\) such that
\[
g(\tilde{t}) = \mathbf{c}.
\]
At time \(\tilde{t}\), the difference between our shot and the submarine’s position is:
\[
f(\tilde{t}) - x(\tilde{t}) = \left( \sum_{j \in \operatorname{dom}(\mathbf{c})} \mathbf{c}(j) \tilde{t}^j \right) - \left( \sum_{i = 0}^{n} c_i \tilde{t}^i \right) = 0.
\]
Thus, the torpedo launched at time \(\tilde{t}\) will hit the submarine.
\\\\
The overall strategy is simply to apply \(f\) recursively:
\[
\text{At time } t = 0, \text{ shoot at } f(0); \quad t = 1, \text{ shoot at } f(1); \quad t = 2, \text{ shoot at } f(2); \quad \dots
\]
By construction, there will be a finite time \(\tilde{t}\) when the torpedo hits the submarine.
\end{solution}
